% GENERATED FILE. Edit canonical lessons, then run npm run generate:books.
% canonical-source-hash: fnv1a64:4c552f95af1b960a
% canonical-lessons: TE-C117-kuragayalu, TE-C117-aratipandu, TE-C117-mamidipandu, TE-C117-kobbari, TE-C117-ullipaya

\chapter{Vegetables, Banana, Mango, Coconut, Onion}
\label{ch:te-vegetables-banana-mango-coconut-onion}
\hlchaptermodality{117}

\begin{chapteropening}
\textbf{By the end of this chapter:} \emph{I can say vegetables, a banana, a mango, coconut and an onion.}

Complete the last lesson of chapter 117: I can say vegetables, a banana, a mango, coconut and an onion.
\end{chapteropening}

\section[\texorpdfstring{\te{కూరగాయలు} (kūragāyalu)}{కూరగాయలు (kūragāyalu)}]{\te{కూరగాయలు} (kūragāyalu) — vegetables}

\label{lesson:TE-C117-kuragayalu}



\begin{quote}
Before the new one: say the Telugu for a fish, then the Telugu for meat.
\end{quote}

\subsection*{You'll want to know: \te{కూరగాయలు}}
\textbf{\te{కూరగాయలు}} — \emph{kūragāyalu} — \textquotedblleft{}vegetables\textquotedblright{}.

\begin{grammarlens}[title={What you've built}]
One more word for this chapter.
\end{grammarlens}

\subsection*{Your turn}
\begin{itemize}
\raggedright
  \item \emph{Say it:} \emph{kūragāyalu}
  \item \emph{Say it:} \emph{kūragāyalu}, once more
  \item \emph{Say it:} say \emph{kūragāyalu}
  \item \emph{Recall:} say the Telugu for a fish, then the Telugu for meat, then say \emph{kūragāyalu} again
\end{itemize}

\begin{tcolorbox}[breakable,colback=teal!4,colframe=teal!35!black,title={Before you move on}]
What does \te{కూరగాయలు} mean? (\textquotedblleft{}Vegetables\textquotedblright{}.) Say it once more.
\end{tcolorbox}
\section[\texorpdfstring{\te{అరటిపండు} (araṭipaṇḍu)}{అరటిపండు (araṭipaṇḍu)}]{\te{అరటిపండు} (araṭipaṇḍu) — a banana}

\label{lesson:TE-C117-aratipandu}



\begin{quote}
Before the new one: say the Telugu for meat, then the Telugu for vegetables.
\end{quote}

\subsection*{You'll want to know: \te{అరటిపండు}}
\textbf{\te{అరటిపండు}} — \emph{araṭipaṇḍu} — \textquotedblleft{}a banana\textquotedblright{}.

\begin{grammarlens}[title={What you've built}]
One more word for this chapter.
\end{grammarlens}

\subsection*{Your turn}
\begin{itemize}
\raggedright
  \item \emph{Say it:} \emph{araṭipaṇḍu}
  \item \emph{Say it:} \emph{araṭipaṇḍu}, once more
  \item \emph{Say it:} say \emph{araṭipaṇḍu}
  \item \emph{Recall:} say the Telugu for meat, then the Telugu for vegetables, then say \emph{araṭipaṇḍu} again
\end{itemize}

\begin{tcolorbox}[breakable,colback=teal!4,colframe=teal!35!black,title={Before you move on}]
What does \te{అరటిపండు} mean? (\textquotedblleft{}A banana\textquotedblright{}.) Say it once more.
\end{tcolorbox}
\section[\texorpdfstring{\te{మామిడిపండు} (māmiḍipaṇḍu)}{మామిడిపండు (māmiḍipaṇḍu)}]{\te{మామిడిపండు} (māmiḍipaṇḍu) — a mango}

\label{lesson:TE-C117-mamidipandu}



\begin{quote}
Before the new one: say the Telugu for vegetables, then the Telugu for a banana.
\end{quote}

\subsection*{You'll want to know: \te{మామిడిపండు}}
\textbf{\te{మామిడిపండు}} — \emph{māmiḍipaṇḍu} — \textquotedblleft{}a mango\textquotedblright{}.

\begin{grammarlens}[title={What you've built}]
One more word for this chapter.
\end{grammarlens}

\subsection*{Your turn}
\begin{itemize}
\raggedright
  \item \emph{Say it:} \emph{māmiḍipaṇḍu}
  \item \emph{Say it:} \emph{māmiḍipaṇḍu}, once more
  \item \emph{Say it:} say \emph{māmiḍipaṇḍu}
  \item \emph{Recall:} say the Telugu for vegetables, then the Telugu for a banana, then say \emph{māmiḍipaṇḍu} again
\end{itemize}

\begin{tcolorbox}[breakable,colback=teal!4,colframe=teal!35!black,title={Before you move on}]
What does \te{మామిడిపండు} mean? (\textquotedblleft{}A mango\textquotedblright{}.) Say it once more.
\end{tcolorbox}
\section[\texorpdfstring{\te{కొబ్బరి} (kobbari)}{కొబ్బరి (kobbari)}]{\te{కొబ్బరి} (kobbari) — coconut}

\label{lesson:TE-C117-kobbari}



\begin{quote}
Before the new one: say the Telugu for a banana, then the Telugu for a mango.
\end{quote}

\subsection*{You'll want to know: \te{కొబ్బరి}}
\textbf{\te{కొబ్బరి}} — \emph{kobbari} — \textquotedblleft{}coconut\textquotedblright{}.

\begin{grammarlens}[title={What you've built}]
One more word for this chapter.
\end{grammarlens}

\subsection*{Your turn}
\begin{itemize}
\raggedright
  \item \emph{Say it:} \emph{kobbari}
  \item \emph{Say it:} \emph{kobbari}, once more
  \item \emph{Say it:} say \emph{kobbari}
  \item \emph{Recall:} say the Telugu for a banana, then the Telugu for a mango, then say \emph{kobbari} again
\end{itemize}

\begin{tcolorbox}[breakable,colback=teal!4,colframe=teal!35!black,title={Before you move on}]
What does \te{కొబ్బరి} mean? (\textquotedblleft{}Coconut\textquotedblright{}.) Say it once more.
\end{tcolorbox}
\section[\texorpdfstring{\te{ఉల్లిపాయ} (ullipāya)}{ఉల్లిపాయ (ullipāya)}]{\te{ఉల్లిపాయ} (ullipāya) — an onion}

\label{lesson:TE-C117-ullipaya}



\begin{quote}
Before the new one: say the Telugu for a mango, then the Telugu for coconut.
\end{quote}

\subsection*{You'll want to know: \te{ఉల్లిపాయ}}
\textbf{\te{ఉల్లిపాయ}} — \emph{ullipāya} — \textquotedblleft{}an onion\textquotedblright{}.

\begin{grammarlens}[title={What you've built}]
One more word for this chapter.
\end{grammarlens}

\subsection*{Your turn}
\begin{itemize}
\raggedright
  \item \emph{Say it:} \emph{ullipāya}
  \item \emph{Say it:} \emph{ullipāya}, once more
  \item \emph{Say it:} say \emph{ullipāya}
  \item \emph{Recall:} say the Telugu for a mango, then the Telugu for coconut, then say \emph{ullipāya} again
\end{itemize}

\begin{tcolorbox}[breakable,colback=teal!4,colframe=teal!35!black,title={Before you move on}]
What does \te{ఉల్లిపాయ} mean? (\textquotedblleft{}An onion\textquotedblright{}.) Say it once more.
\end{tcolorbox}
